\subsection{Solana}

Solana claims to solve the trilemma:

\bquote{
    In this post, we'll explore Turbine, Solana's block propagation protocol---inspired by BitTorrent---that solves the blockchain scalability trilemma.
    % [\dots] \\
    % With that context, let's jump into Turbine, Solana's block propagation protocol, to explain how the Solana network propagates data to solve the blockchain scalability trilemma.
}[\linkSolanaTurbineTrilemma (Anatoly Yakovenko, 2019-07-11)]

In March 2021, Anatoly Yakovenko --- author of the Solana whitepaper and a co-founder --- said this in answer to the question ``How does Solana solve the trilemma: decentralized, scalable and secure?'':

\bquote{
    With hardware. Latest nvidia GPU has 10,000 cores. So imagine each validator doing 10,000 more work [than] a single threaded evm environment.
}[\linkSolanaRedditTrilemma (Anatoly Yakovenko, 2021-03-13)]

These quotes contradict one another.

%
%

\subsubsection{Block Propagation?}

\label{sec:solana-block-propagation}

In \linkSolanaTurbineTrilemma, Yakovenko says:

\bquote{
    The scalability trilemma in blockchain technology is all about bandwidth. In most blockchain networks today, given a fixed amount of bandwidth per node, increasing the node count will increase the amount of time necessary to propagate all the data to all nodes. That's a big problem.
}

This is inaccurate --- the scalability criterion of the trilemma is about \emph{capacity to process transactions}.

Reliable block propagation is important for blockchains and \emph{can} be a capacity bottleneck.
However, good block propagation doesn't solve the trilemma --- just because block propagation isn't a bottleneck does not mean that a network can scale.

This misstatement of the trilemma constitutes changing the goal posts.

\subsubsection{GPUs are still $O(c)$}

Yakovenko's other idea (from March 2021) is to optimize each node's ability to process transactions via consumer hardware.
But, that consumer hardware (a GPU) is an $O(c)$ device --- adding a GPU to a computer does not allow that single computer to process $> O(c)$ transactions.
In this case, the bottleneck is still an $O(c)$ device, so the network is still limited to $O(c)$ scalability.

% What is ``that context''?
% Under the section ``The Scalability Trilemma'':

% \bquote{
%     The scalability trilemma in blockchain technology is all about bandwidth. [\dots]
% }

% \bquote{
%     With hardware. Latest nvidia GPU has 10,000 cores. So imagine each validator doing 10,000 more work then a single threaded evm environment.
% }[\linkSolanaRedditTrilemma]


% not solana, does not come with a citation:

% \bquote{
%     Solana relies on proof of history as the key innovation to solve the blockchain trilemma instead of sharding or relying on Layer-2s.
% }[Martin Lee; \href{https://www.nansen.ai/research/solana-scalability-through-speed}{Solana: Scalability through speed}]
